% Nguồn SGK Toán 9 Tập một — Chương III, Luyện tập chung sau Bài 10, trang 63--64:
% https://cdn3.olm.vn/upload/thuvienso/4491668113-page-63-1773022941473.png
% https://cdn3.olm.vn/upload/thuvienso/4491668113-page-64-1773022941797.png

\section*{Luyện tập chung}

\begin{lesson}
    \textbf{Ví dụ 1}

    Rút gọn biểu thức
    \[
        A=\sqrt{(\sqrt{3}+1)^2}-\sqrt{3(\sqrt{3}+1)^2}+\sqrt{4}.
    \]

    \textbf{Giải}

    Ta có
    \[
        A=\sqrt{(\sqrt{3}+1)^2}-\sqrt{3}\cdot\sqrt{(\sqrt{3}+1)^2}+\sqrt{2^2}
        =(\sqrt{3}+1)-\sqrt{3}(\sqrt{3}+1)+2=0.
    \]

    \textbf{Ví dụ 2}

    Cho biểu thức
    \[
        P=\dfrac{\sqrt{x}}{\sqrt{x}+1}+\dfrac{\sqrt{x}}{\sqrt{x}-1}
    \]
    với $x>1$.

    a) Rút gọn $P$.

    b) Sử dụng kết quả câu a, tính giá trị của $P$ khi $x=101$.

    \textbf{Giải}

    a)
    \[
        P=\dfrac{\sqrt{x}(\sqrt{x}-1)+\sqrt{x}(\sqrt{x}+1)}{(\sqrt{x}+1)(\sqrt{x}-1)}
        =\dfrac{x-\sqrt{x}+x+\sqrt{x}}{x-1}
        =\dfrac{2x}{x-1}.
    \]

    b) Khi $x=101$ thì
    \[
        P=\dfrac{2\cdot101}{101-1}=\dfrac{202}{100}=2,02.
    \]

    \textbf{Ví dụ 3}

    Người ta cần làm một thùng hình lập phương bằng bìa cứng không có nắp trên và có thể tích $216$ dm$^3$ để đựng đồ. Tính diện tích bìa cứng cần dùng để làm thùng đựng đó (coi diện tích các mép nối là không đáng kể).

    % TODO(HINH-NGUON): bo sung asset/crop dung hinh thung lap phuong khong nap tu SGK trang 63; khong redraw phoi canh xap xi bang TikZ.

    \textbf{Giải}

    Gọi $x$ (dm, $x>0$) là độ dài cạnh của thùng hình lập phương cần làm.

    Ta có: $x^3=216$, suy ra $x=\sqrt[3]{216}=6$ (dm).

    Vì thùng đựng không có nắp nên thùng gồm $4$ mặt bên và $1$ mặt đáy, mỗi mặt là một hình vuông cạnh $6$ dm. Do đó diện tích bìa cứng cần dùng là
    \[
        S=5\cdot6^2=180\,(\text{dm}^2).
    \]
\end{lesson}

\begin{exercises}
    \subsection{Bài tập}

    \begin{questions}[resume=SGK]
        \item Rút gọn các biểu thức sau:
        \begin{questions}
            \item $\dfrac{5+3\sqrt{5}}{\sqrt{5}}-\dfrac{1}{\sqrt{5}-2}$;

            \answerline
            \answerline

            \item $\sqrt{(\sqrt{7}-2)^2}-\sqrt{63}+\dfrac{\sqrt{56}}{\sqrt{2}}$;

            \answerline
            \answerline

            \item $\dfrac{\sqrt{(\sqrt{3}+\sqrt{2})^2}+\sqrt{(\sqrt{3}-\sqrt{2})^2}}{2\sqrt{12}}$;

            \answerline
            \answerline

            \item $\dfrac{\sqrt[3]{(\sqrt{2}+1)^3}-1}{\sqrt{50}}$.

            \answerline
            \answerline
        \end{questions}

        \item Tính giá trị của các biểu thức sau:
        \begin{questions}
            \item $3\sqrt{45}+\dfrac{5\sqrt{15}}{\sqrt{3}}-2\sqrt{245}$;

            \answerline
            \answerline

            \item $\dfrac{\sqrt{12}-\sqrt{4}}{\sqrt{3}-1}-\dfrac{\sqrt{21}+\sqrt{7}}{\sqrt{3}+1}+\sqrt{7}$;

            \answerline
            \answerline

            \item $\dfrac{3-\sqrt{3}}{1-\sqrt{3}}+\sqrt{3}(2\sqrt{3}-1)+\sqrt{12}$;

            \answerline
            \answerline

            \item $\dfrac{\sqrt{3}-1}{\sqrt{2}}+\dfrac{\sqrt{2}}{\sqrt{3}-1}-\dfrac{6}{\sqrt{6}}$.

            \answerline
            \answerline
        \end{questions}

        \item Giả sử lực $F$ của gió khi thổi theo phương vuông góc với bề mặt cánh buồm của một con thuyền tỉ lệ thuận với bình phương tốc độ của gió, hệ số tỉ lệ là $30$. Trong đó, lực $F$ được tính bằng N (Newton) và tốc độ được tính bằng m/s.
        \begin{questions}
            \item Khi tốc độ của gió là $10$ m/s thì lực $F$ là bao nhiêu Newton?

            \answerline
            \answerline

            \item Nếu cánh buồm chỉ có thể chịu được một áp lực tối đa là $12\,000$ N thì con thuyền đó có thể đi được trong gió với tốc độ gió tối đa là bao nhiêu?

            \answerline
            \answerline
        \end{questions}

        \item Rút gọn các biểu thức sau:
        \begin{questions}
            \item $\sqrt[3]{(-x-1)^3}$;

            \answerline
            \answerline

            \item $\sqrt[3]{8x^3-12x^2+6x-1}$.

            \answerline
            \answerline
        \end{questions}
    \end{questions}
\end{exercises}
